\section{Introducción}\label{sec:intro}

El reto AgroCebada de FIRA 2026 pide estimar el rendimiento de cebada de temporal, en toneladas por hectárea (t/ha),
de 59 parcelas de Hidalgo, Puebla y Tlaxcala. Estas tres entidades se cuentan entre las cuatro principales productoras
de cebada grano del país~\cite{sader2025}. Para entrenar tenemos 138 parcelas del mismo ciclo de 2025 con su
rendimiento declarado, imágenes de los satélites Sentinel-2~\cite{drusch2012} y PlanetScope~\cite{frazier2021},
clima~\cite{funk2015,verdin2020} y topografía. El jurado califica primero el error de las 59 predicciones con el RMSE,
que se mide en t/ha y castiga más los errores grandes.

La predicción de rendimiento con aprendizaje automático y teledetección es un tema muy estudiado~\cite{vanklompenburg2020,weiss2020}.
En este reto las parcelas son pequeñas y su rendimiento es declarado, no medido, y la literatura reporta que ese tipo
de dato contiene error considerable~\cite{lobell2020}. Por eso el reporte se centra en qué información aporta de
verdad a la predicción y en qué tan confiable es el error que estimamos.

Este reporte explica cómo construimos el modelo y por qué tomamos cada decisión. La Fig.~\ref{fig:flujo} resume el
proceso. Partimos de un hallazgo en las etiquetas (Sección~\ref{sec:datos}), fijamos cómo íbamos a evaluar cada idea
(Sección~\ref{sec:evaluacion}) y agregamos piezas una por una, conservando solo las que mejoraban la validación
(Sección~\ref{sec:construccion}). Después describimos el modelo final (Sección~\ref{sec:final}), qué tan confiable es
su error (Sección~\ref{sec:confiabilidad}) y qué aportan los índices satelitales (Sección~\ref{sec:indices}).

\begin{figure*}[t]
\centering
\begin{tikzpicture}[node distance=0.45cm, every node/.style={font=\footnotesize, align=center},
  caja/.style={draw=black!60, rounded corners=2pt, minimum height=1.05cm, text width=2.35cm, fill=black!3},
  clave/.style={caja, draw=blue!60!black, fill=blue!6}, flecha/.style={-{Stealth[length=2mm]}, black!60}]
\node[caja] (d) {Datos\\138 + 59 parcelas\\satélite, clima, suelo};
\node[caja, right=of d] (v) {91 variables\\por parcela};
\node[clave, right=of v] (r) {Régimen de la\\etiqueta (A o B)};
\node[clave, right=of r] (z) {Rezago espacial\\vecinos a 400\,m};
\node[clave, right=of z] (g) {Proceso\\gaussiano};
\node[caja, right=of g] (s) {Predicción\\e intervalo\\p10 a p90};
\draw[flecha] (d) -- (v); \draw[flecha] (v) -- (r); \draw[flecha] (r) -- (z); \draw[flecha] (z) -- (g); \draw[flecha] (g) -- (s);
\node[below=0.25cm of v, text width=2.6cm, font=\scriptsize, text=black!60] {probadas, sin mejora\\(Sección~\ref{sec:bloques})};
\end{tikzpicture}
\caption{Arquitectura del modelo. En azul, las piezas que usa el modelo final. Las 91 variables satelitales, de clima
y de topografía se probaron una por bloque y ninguna mejoró la validación.}
\label{fig:flujo}
\end{figure*}
